\documentclass[
11pt, 
a4paper,
DIV = 14,
twoside,
twocolumn, % z w e i s e i t i g e r A u f b a u
parskip =half, % H a l b e Z e i l e A b s a t z a b s t a n d
headsepline, % K o p f z e i l e n s e p a r a t i o n s l i n i e
openright, % N e u e s K a p i t e l b e g i n n t r e c h t s
]{scrreprt}

\usepackage{lipsum}
\usepackage{graphicx}
%\usepackage[chapter]{placeins}
\usepackage{fontenc}
\usepackage{scrlayer-scrpage}
\usepackage[ngerman]{babel}
\usepackage{lmodern}
\usepackage{blindtext}
\usepackage[utf8]{inputenc}
\usepackage[backend=biber,style=alphabetic]{biblatex}
\usepackage{csquotes}
\addbibresource{ref/ref.bib}


\begin{document}
	\ohead{\includegraphics[height=45pt]{bilder/thi.png.png}}
	
	%%%%%%%%% KOPF %%%%%%%%%%%%%%%%%%%%%%%%%%
	\title{Welche Rolle spielen künstliche Intelligenz und Machine Learning bei der Optimierung von Energieversorgungssystemen?}
	\author{Nico Elsner\\
		Studiengang: Robotik}
	\date{\today}
	\maketitle

	% Seitenzahl ausblenden bis nach Inhaltsverzeichnis
	\clearpage
	\begingroup
	  \pagestyle{empty}
	  \renewcommand*\chapterpagestyle{empty}
	  \tableofcontents 
	\clearpage
	\endgroup
	
	\thispagestyle{empty}

	%% Einleitung
	\chapter{Einleitung}
	\setcounter{page}{1}
	\section{Hintergrund und Motivation des Themas}
	Hintergrund dieser Arbeit ist die in letzter Zeit extrem gestiegene Die
Energieversorgungssysteme stehen heutzutage vor einer Vielzahl von
Herausforderungen. Der steigende Energiebedarf, die begrenzte Verfügbarkeit
konventioneller Energieressourcen und die zunehmende Bedeutung erneuerbarer
Energien erfordern innovative Lösungen, um die Effizienz und Nachhaltigkeit der
Energieversorgung zu verbessern. In diesem Zusammenhang spielen künstliche
Intelligenz (KI) und maschinelles Lernen (ML) eine immer wichtigere Rolle.

Die Grundidee hinter dem Einsatz von KI und ML in der Optimierung von
Energieversorgungssystemen besteht darin, dass diese Technologien in der Lage
sind, große Mengen an Daten zu analysieren, Muster zu erkennen, Vorhersagen zu
treffen und adaptive Entscheidungen zu treffen. Dadurch können
Energieversorgungssysteme effizienter gestaltet, die Zuverlässigkeit erhöht und
Kosten gesenkt werden.

Ein zentrales Anwendungsgebiet von KI und ML liegt in der Prognose von
Energieerzeugung und -nachfrage. Durch den Einsatz von ML-Algorithmen können
präzisere Vorhersagen über den zukünftigen Energiebedarf gemacht werden, was
eine optimierte Planung der Energieerzeugung und -verteilung ermöglicht.
Darüber hinaus können ML-Modelle genutzt werden, um Lastmanagement-Strategien
zu entwickeln, um den Energieverbrauch zu optimieren und die Spitzenlasten zu
reduzieren.

Ein weiteres wichtiges Einsatzgebiet von KI und ML liegt in der Optimierung von
Stromnetzen und -verteilung. Hier können intelligente Algorithmen eingesetzt
werden, um den Energiefluss zu steuern, Engpässe zu identifizieren und die
Netzstabilität zu gewährleisten. ML kann auch bei der Erkennung von Fehlern und
Störungen in Energieversorgungssystemen helfen, indem Anomalien in den Daten
erkannt und potenzielle Probleme frühzeitig gemeldet werden.

Die Integration erneuerbarer Energien stellt ebenfalls eine große
Herausforderung dar, da sie von Natur aus volatil und schwer vorhersehbar sind.
Hier kommen KI und ML zum Einsatz, um Vorhersagemodelle für die Stromerzeugung
aus erneuerbaren Quellen zu entwickeln, um die Integration in das Gesamtsystem
zu optimieren. Dies ermöglicht eine bessere Ausnutzung erneuerbarer Energien
und eine verbesserte Integration mit konventionellen Energiequellen.

Neben den technischen Aspekten gibt es auch weitere wichtige Faktoren, die bei
der Optimierung von Energieversorgungssystemen berücksichtigt werden müssen.
Datenschutz und Sicherheit spielen eine entscheidende Rolle, da die
Verarbeitung großer Mengen an Energie- und Verbrauchsdaten eine Herausforderung
für den Schutz sensibler Informationen darstellt. Darüber hinaus stellen sich
auch ethische Fragen, wie beispielsweise der faire Zugang zu Energie oder die
Vermeidung von Diskriminierung bei der Nutzung von KI-gesteuerten
Energiesystemen.

In Anbetracht all dieser Aspekte ist es von großer Bedeutung, die Rolle von KI
und ML bei der Optimierung von Energieversorgungssystemen zu untersuchen. Durch
die Nutzung des Potenzials dieser Technologien können wir eine effizientere,
nachhaltigere und zuverlässigere Energieversorgung erreichen. In dieser
Seminararbeit werden wir die verschiedenen Einsatzbereiche von KI und ML in der
Energieversorgung untersuchen und ihre Auswirkungen sowie mögliche
Herausforderungen und zukünftige Entwicklungen analysieren.
	%\section{Zielsetzung und Fragestellung der Arbeit}
	%\ldots


	\chapter{Grundlagen der künstlichen Intelligenz und Machine Learning}
	\input{kapitel/Grundlagen der künstlichen Intelligenz und Machine Learning.tex}
	%\section{Definition und Abgrenzung}
	%\ldots
	%\section{Typen von Machine Learning Algorithmen}
	%\ldots


	\chapter{Optimierung von Energieversorgungssystemen}
	\input{kapitel/Optimierung von Energieversorgungssystemen.tex}
	\section{Überblick über Energieversorgungssystemen}
	\input{kapitel/Überblick von Energieversorgungssystemen.tex}
	\section{Herausforderungen bei der Optimierung}
	Die Optimierung von Energieversorgungssystemen steht vor einer Reihe von
Herausforderungen, die es zu bewältigen gilt. Eine der zentralen
Herausforderungen besteht in der Komplexität dieser Systeme.
Energieversorgungssysteme bestehen aus einer Vielzahl von Komponenten und
Akteuren, wie beispielsweise Energieerzeugung, -übertragung, -verteilung und
-verbrauch. Die Koordination und Optimierung all dieser Komponenten erfordert
eine umfassende Analyse und Steuerung, um effiziente Ergebnisse zu erzielen.

Ein weiterer wichtiger Aspekt ist die Volatilität erneuerbarer Energien. Die
Integration von erneuerbaren Energien wie Solarenergie und Windenergie in das
Energieversorgungssystem stellt eine Herausforderung dar. Diese Energiequellen
unterliegen starken Schwankungen aufgrund von Wetterbedingungen und weisen
volatilere Erzeugungsmuster auf als konventionelle Energieressourcen.

Darüber hinaus stehen Energieversorgungssysteme vor Herausforderungen im
Zusammenhang mit der Skalierbarkeit und Flexibilität. Da die Energienachfrage
schwankt und sich die Energieerzeugungsmuster ändern, müssen die Systeme in der
Lage sein, sich an diese Veränderungen anzupassen und eine zuverlässige
Energieversorgung sicherzustellen.

Die Sicherheit und Zuverlässigkeit des Energieversorgungssystems ist eine
weitere Herausforderung. Da das Energienetz ein attraktives Ziel für
Cyberangriffe darstellt, müssen Maßnahmen ergriffen werden, um die Systeme vor
potenziellen Bedrohungen zu schützen und die Verfügbarkeit der
Energieversorgung zu gewährleisten.

Nicht zuletzt sind auch wirtschaftliche und regulatorische Herausforderungen zu
berücksichtigen. Die Optimierung von Energieversorgungssystemen erfordert
Investitionen in Infrastruktur, Technologien und Forschung. Zudem müssen
regulatorische Rahmenbedingungen geschaffen werden, die die Integration neuer
Technologien und Geschäftsmodelle ermöglichen und den Übergang zu einer
nachhaltigen Energieversorgung unterstützen.

Insgesamt stehen die Optimierung von Energieversorgungssystemen vor
verschiedenen Herausforderungen, die von der Komplexität der Systeme über die
Integration erneuerbarer Energien bis hin zu Fragen der Sicherheit,
Skalierbarkeit und Wirtschaftlichkeit reichen. Die Bewältigung dieser
Herausforderungen erfordert einen interdisziplinären Ansatz, der
technologische, wirtschaftliche, regulatorische und soziale Aspekte
berücksichtigt.

	\section{Allgemeine Möglichkeiten der Optimierung}
	\input{kapitel/Allgemeine Möglichkeiten der Optimierung.tex}
	

	\chapter{Rolle von künstlicher Intelligenz und Machine Learning bei der Optimierung von Energieversorgungssystemen}
	\input{kapitel/Rolle von künstlicher Intelligenz und Machine Learning bei der Optimierung von Energieversorgungssystemen.tex}
	\section{Potenzial und Vorteile von künstlicher Intelligenz und Machine Learning}
	\input{kapitel/Potenzial und Vorteile von künstlicher Intelligenz und Machine Learning.tex}
	\section{Beispiele und Anwendungen von künstlicher Intelligenz und Machine Learning in der Energieversorgung}
	\input{kapitel/Beispiele und Anwendungen von künstlicher Intelligenz und Machine Learning in der Energieversorgung.tex}


	\chapter{Kritische Betrachtung und Ausblick}
	\input{kapitel/Kritische Betrachtung und Ausblick.tex}
	\section{Grenzen und Risiken von künstlicher Intelligenz}
	\input{kapitel/Grenzen und Risiken von künstlicher Intelligenz.tex}
	\section{Perspektiven}
	Die Rolle von künstlicher Intelligenz und maschinellem Lernen bei der
Optimierung von Energieversorgungssystemen bietet vielversprechende
Perspektiven für die Zukunft. Das die Softwareindustrie, welche die KI-Systeme
maßgeblich mitprägt, immer schneller am wachsen ist zeigt
Abbildung(\ref{fig:wachstumsberufe}). Unter allen Berufen, ist der Informatiker
mit großem Abstand der schnellst Wachsende.
\begin{figure}[!h]
    \centering
    \includesvg[width=1\columnwidth]{bilder/Wachstumsberufe Deutschland_2.svg}
    \caption[width=0.4\textwidth]{Ablauf bei Reinforcement Learning}
    \label{fig:wachstumsberufe}
\end{figure}
\cite{eichhorst2015zukunft}

Mit dem weiteren Fortschritt von KI- und ML-Technologien eröffnen sich neue
Möglichkeiten zur Effizienzsteigerung, Kostensenkung und Verbesserung der
Nachhaltigkeit in der Energieversorgung.

Eine vielversprechende Perspektive besteht in der Integration von KI und ML in
Smart Grids und intelligenten Energienetzen. Durch den Einsatz
fortschrittlicher Algorithmen und Modelle können Energieflüsse optimiert,
Spitzenlasten reduziert und erneuerbare Energien effizienter in das
Energiesystem integriert werden. Dies ermöglicht eine bessere Ausbalancierung
von Angebot und Nachfrage und trägt zu einer stabilen und nachhaltigen
Energieversorgung bei.

Des Weiteren eröffnet die Kombination von KI und ML mit dem Internet der Dinge
(IoT) und dem Einsatz von Sensoren neue Perspektiven für die Energieeffizienz.
Intelligente Geräte und Systeme können Energieverbrauchsdaten erfassen und
analysieren, um personalisierte Empfehlungen zur Reduzierung des
Energieverbrauchs zu geben. Dies ermöglicht es Verbrauchern, bewusstere
Entscheidungen zu treffen und ihren Energieverbrauch zu optimieren.

Ein weiteres vielversprechendes Feld ist die Entwicklung autonomer
Energieversorgungssysteme. Durch den Einsatz von KI und ML können autonome
Energieerzeugungs- und Speicherlösungen entwickelt werden, die den Bedarf an
menschlichem Eingriff minimieren. Diese Systeme können den Energieverbrauch und
die -erzeugung in Echtzeit überwachen, Vorhersagen treffen und automatisch
optimale Entscheidungen treffen, um die Effizienz und Zuverlässigkeit der
Energieversorgung zu maximieren.

Darüber hinaus bieten KI und ML Möglichkeiten zur besseren Vorhersage und
Bewältigung von Energienotfällen. Durch die Analyse von Echtzeitdaten können
frühzeitige Warnungen und reaktionsschnelle Maßnahmen bei Störungen im
Energiesystem ermöglicht werden. Dies unterstützt die Sicherheit und Resilienz
der Energieinfrastruktur.

Insgesamt eröffnen die Fortschritte in KI und ML spannende Perspektiven für die
Energieversorgungssysteme der Zukunft. Mit einer kontinuierlichen Forschung und
Entwicklung, der Verfügbarkeit hochwertiger Daten und einer sorgfältigen
Integration dieser Technologien können wir eine effizientere, nachhaltigere und
zuverlässigere Energieversorgung erreichen.


	\chapter{Fazit und Ausblick}
	Die Rolle von künstlicher Intelligenz (KI) und maschinellem Lernen (ML) bei der
Optimierung von Energieversorgungssystemen ist von großer Bedeutung. Die
Anwendung von KI und ML bietet vielfältige Möglichkeiten zur
Effizienzsteigerung, Kostenoptimierung, Integration erneuerbarer Energien und
Verbesserung des Verbraucherverhaltens. Durch den Einsatz von KI können
Energieversorgungsunternehmen fundierte Entscheidungen treffen, Prognosen
erstellen, den Energiefluss optimieren und die Zuverlässigkeit der
Energieversorgung verbessern.

Jedoch sind bei der Nutzung von KI in der Energieversorgung auch
Herausforderungen zu berücksichtigen. Grenzen wie die Qualität und
Verfügbarkeit von Daten, ethische Überlegungen, potenzielle
Fehlinterpretationen oder Fehlentscheidungen von KI-Modellen sowie
Sicherheitsrisiken erfordern eine sorgfältige Planung, Überwachung und
Implementierung.

Ausblick: Für die Zukunft bieten sich auf dem Gebiet der künstlichen
Intelligenz und des maschinellen Lernens in der Energieversorgung zahlreiche
Möglichkeiten. Fortschritte in der Datenverfügbarkeit und -qualität werden die
Genauigkeit und Leistungsfähigkeit von KI-Modellen verbessern. Eine verstärkte
Berücksichtigung ethischer Prinzipien bei der Entwicklung und Implementierung
von KI-Systemen wird zu einer fairen und transparenten Nutzung beitragen.

Darüber hinaus könnten Fortschritte im Bereich des Verständnis von
KI-Modellen dazu beitragen, Vertrauen und Akzeptanz in deren Entscheidungen zu
fördern. Die Integration von KI-Systemen mit dem Internet der Dinge (IoT) und
anderen fortschrittlichen Technologien ermöglicht eine noch intelligentere und
effizientere Energieversorgung.

Ein wichtiger Aspekt für den Ausblick ist die kontinuierliche Forschung und
Entwicklung auf dem Gebiet der künstlichen Intelligenz und des maschinellen
Lernens in der Energieversorgung. Neue Algorithmen, Modelle und Techniken
werden entwickelt, um die Herausforderungen anzugehen und die
Leistungsfähigkeit von KI-Systemen weiter zu verbessern.

Insgesamt bietet die Kombination von künstlicher Intelligenz und maschinellem
Lernen enorme Potenziale für die Optimierung von Energieversorgungssystemen.
Mit einer sorgfältigen Betrachtung der Grenzen, Risiken und Chancen können wir
eine nachhaltigere, effizientere und zuverlässigere Energieversorgung der
Zukunft erreichen.

\end{document}